Paper~F~\cite{F} proved the logarithmic law of the resolution gap with the sharp constant for periodic windows with at most twenty harmonics or with decay exponent $K\ge0.003$ (which includes every window with $x\ge0.003$):
with $\ell=\ln(1/\sigma)$, the smallest gap obeys explicit two-sided bounds on $g_{\min}\sqrt{\ell}$ for $10^{-100000}\le\sigma\le10^{-24}$, and its limit lies between $1-\bar\varepsilon$ and $1$, where $\bar\varepsilon$ bounds the remainder on the window class.
Here we extend the lower side to windows with any number of harmonics whose ripple is weak in the sup norm or poor in low-frequency energy.
Three tools do the work.
A sharper treatment of one remainder of the curvature step turns an absolute error of size $e^{-\tau/2}/k'$ into a relative one.
A split sandwich lets a low-frequency block enter through its energy rather than through its sup.
Cell-based computer-assisted proofs, with an explicit bound on floating-point rounding in every acceptance test, certify the full lower bound, remainders included, and a mean-square bound of the curvature cross term closes the edge of very wide focus.
The main theorem gives $g\sqrt{\ell}\ge(1-\bar\varepsilon)(\lambda/\Upsilon)(1+12.1/\ell+2.5\times10^{-4})$ for every state with $\tau\ge20$, for $\ell\in[55.26,960.8]$, that is, for $\sigma$ from about $10^{-24}$ down to about $10^{-417}$.
In this range the coefficient one holds, up to the factor $1-\bar\varepsilon$, for every window with $\eta\le1/\ell$, for a family of intermediate windows, and for seventeen thin-block windows certified one by one.
A conjecture of our own (that $\eta e^{\mu/2}$ stays bounded) is false; we give explicit counterexamples, and they are covered by the theorem.
We close with the road to an unconditional theorem, from the nearest restriction to the farthest.
For the states with a narrow focus ($\tau<20$) we prove that positivity alone makes the local brightness pay for fineness, that a ripple hidden at the centre of the focus is dark, and that the local route costs almost no remainder; with these bounds the envelope exceeds the target at every $\ell$ computed, but the cells that would turn this into a proof are not yet done.
Further out lie the full bound beyond $\ell=960.8$ and the limit $\sigma\to0$, windows outside the covered families, the class constant $\bar\varepsilon$, and the upper side.
We list exactly what is proved, what is computer-assisted, and what is only numerical.
This version (v1.3) states the hypothesis of the lemma for very wide focus for the full range $\ell\ge55.26$ (up to v1.2 it was stated for $\sigma\le10^{-24}$, which leaves out $\ell\in[55.26,24\ln10)$), makes a constant of the local route a proved bound, and corrects three smaller points of the open problems and of the code; the statement and the range of the main theorem are unchanged.
Version~1.2 corrected two steps of the proof of the main theorem, the lower end point in the brightness $G$ and the step that restores the sandwich factors, and closed both by re-running the certificates; the statement and the range of the theorem were unchanged.
Version~1.1 closed three gaps in the covering of the computer-assisted proofs of v1.0 (a wedge of states inside each column of the main certificate, the extreme values of the energy frequency $\mu$, and the smallest $b$ of the reduced certificate), added a line-by-line covering argument with a random audit of $10^5$ states, and corrected two statements of v1.0 about the narrow-focus states.
